% Part: proof-theory
% Chapter: natural-deduction
% Section: grafting

\documentclass[../../../include/open-logic-section]{subfiles}

\begin{document}

\olfileid{pt}{nat}{gra}

\olsection{\usetoken{P}{proof}નું કલમ-જોડાણ}

\Log{N1c} અને~\Log{N1i}માંની
!!^{proof}s એ ધારણાઓ~$!B$માંથી
!!{formula}s~$!A$ની !!{proof}s છે.
માનો કે~$!B$ની પણ !!a{proof} છે.
તો~$!B$ની !!{proof} અને
$!B$માંથી~$!A$ની !!{proof}
જોડીને~$!A$ની એવી
!!a{proof} કેવી રીતે
મેળવવી જેમાં~$!B$ની ધારણાનો
ઉપયોગ ન થતો હોય?

એક વિકલ્પ એ છે કે~\Intro{\lif}
લગાવી~$!A$ની !!{proof}માંથી
ધારણા~$!B$ દૂર કરીએ.
પછી મળેલી~$!B \lif !A$ની
!!{proof}ને~$!B$ની !!{proof}
સાથે જોડીને આ મળે:
\[
\AxiomC{$\Discharge{!B}{x}$}
\DeduceC{$!A$}
\DischargeRule{\Intro{\lif}}{x}
\UnaryInfC{$!B \lif !A$}
\AxiomC{}
\DeduceC{$!B$}
\RightLabel{\Elim{\lif}}
\BinaryInfC{$!A$}
\DisplayProof
\]
\sourcecorrection{OLMD-224}{મૂળ વૃક્ષમાં અનુસરણ-પ્રવેશનું નિષ્કર્ષ $!B \lif !C$ લખાયું છે; ઉપરની ઉપસાબિતીનું નિષ્કર્ષ $!A$ હોવાથી અહીં $!B \lif !A$ જોઈએ.}
આ રીત સીધી છે, પણ થોડી
અકુદરતી લાગે છે: તરત જ
દૂર કરવાના~$\lif$ને શા માટે
પહેલાં દાખલ કરવું?
$!B \lif !A$માંથી પસાર
થતો આવો ``ચકરાવો''
ટાળી શકાય એવો હોવો જોઈએ.
(આવા ચકરાવા હંમેશાં ટાળી
શકાય છે એ જ સામાન્યીકરણ
પ્રમેયનો આશય છે.)

\Log{N1c} અને~\Log{N1i}માં
આ કિસ્સામાં ચકરાવો સહેલાઈથી
ટાળી શકાય છે: ધારણાઓ~$!B^x$
જ્યાં આવે ત્યાં~$!B$ની
આખી !!{proof} મૂકી દઈએ.
બીજી સાબિતીની ધારણાઓ પર
!!{proof}sનું આવું
``કલમ-જોડાણ'' ઉપયોગી છે,
એટલે તેને વ્યાખ્યા તરીકે
નોંધીએ છીએ.

\begin{defn}
જો~$\delta_1$ એ ખુલ્લી
ધારણા~$!B^x$માંથી~$!A$ની
\Log{N1c}-!!{proof}
(અથવા~\Log{N1i}-!!{proof})
હોય, અને~$\delta_2$ એ~$!B$ની
\Log{N1c}-!!{proof}
(\Log{N1i}-!!{proof}) હોય,
તો~$\delta_1$માંનાં
મુક્ત ન કરાયેલાં~$!B^x$નાં
દરેક સ્થાનને~$\delta_2$થી
બદલવાથી મળતું પરિણામ
$\Subst{\delta_1}{\delta_2}{!B^x}$
છે.
\end{defn}

જો~$\delta_1$ અને~$\delta_2$માં
જુદાં ધારણા-!!{formula}sને
એક જ ચિહ્ન ન આપેલું હોય,
અને~$\delta_2$ની કોઈ ખુલ્લી
ધારણામાં~$\delta_1$નો
સ્વનિર્ધારક ચલ ન આવતો હોય,
તો મળેલી
$\Subst{\delta_1}{\delta_2}{!B^x}$
યોગ્ય !!{proof} છે.
તેની ખુલ્લી ધારણાઓ
$\delta_1$ તથા~$\delta_2$ની
ખુલ્લી ધારણાઓ છે.
!!{proof}sનું કલમ-જોડાણ
કરતી વખતે આ શરતો સામાન્ય રીતે
પૂરી થાય છે. $\delta_1$ અને
$\delta_2$ કોઈ મોટી
!!{proof}ની ઉપસાબિતીઓ હોય
તો પહેલી શરત પૂરી થાય છે.
બીજી શરત પૂરી ન થતી હોય
તો~$\delta_1$ના સ્વનિર્ધારક
ચલોનાં નામ બદલીને તેને
પૂરી કરી શકાય છે.

\Log{N2c} અને~\Log{N2i}
માટે પણ અનુરૂપ વ્યાખ્યા છે.

\begin{defn}
જો~$\delta_1$ એ
$x:!B, \Gamma_1 \Sequent !A$ની
\Log{N2c}-!!{proof}
(અથવા~\Log{N2i}-!!{proof})
હોય, અને~$\delta_2$ એ
$\Gamma_2 \Sequent !B$ની
\Log{N2c}-!!{proof}
(\Log{N2i}-!!{proof}) હોય,
તો~$\delta_1$માંનાં
દરેક સ્વયંસિદ્ધ
$x:!B \Sequent !B$ને
$\delta_2$થી બદલીને,
અને એવાં સ્વયંસિદ્ધોની
નીચે આવતા દરેક સિક્વન્ટના
સંદર્ભમાં~$\Gamma_2$ ઉમેરવાથી
મળતું પરિણામ
$\Subst{\delta_1}{\delta_2}{x:!B}$
છે.
\sourcecorrection{OLMD-225}{મૂળની આ N2 વ્યાખ્યામાં $\delta_2$ને ભૂલથી N1c/N1i-સાબિતી કહેવાયું છે; સિક્વન્ટવાળી આધારસાબિતી માટે અનુરૂપ N2c/N2i જોઈએ.}
\end{defn}

\begin{prop}
જો $\delta_1 \Proves x:!B, \Gamma_1 \Sequent
!A$ અને $\delta_2 \Proves \Gamma_2 \Sequent !B$,
તો, જો~$\delta_1$નો કોઈ
સ્વનિર્ધારક ચલ~$\Gamma_2$માં
ન આવતો હોય,
$\Subst{\delta_1}{\delta_2}{x:!B} \Proves \Gamma_1, \Gamma_2 \Sequent
!A$.
\end{prop}


\begin{proof}
  $\pheight{\delta_1}$ પર આગમનથી.
  \sourcecorrection{OLMD-226}{મૂળની એક-વાક્યીય સાબિતીમાં અપરિભાષિત $\delta$ની ઊંચાઈ લખાઈ છે; પ્રતિજ્ઞાપનમાં આવેલી બદલાતી સાબિતી $\delta_1$ પર આગમનનો આશય છે.}
\end{proof}

\end{document}
